% ---------------------------------------------------------------------------
% Preámbulo compartido (v2): tipografía, color, cuadros, figuras y cabeceras.
% Compilar con XeLaTeX. Todo lo que depende de papers.json llega por
% build/metadata.tex como macros semánticas (\papertitle, \paperauthors, ...).
% ---------------------------------------------------------------------------
\usepackage{fontspec}
% Serifa para el cuerpo: TeX Gyre Pagella (incluida en TeX Live completo).
\setmainfont{texgyrepagella}[Extension=.otf,UprightFont=*-regular,BoldFont=*-bold,
  ItalicFont=*-italic,BoldItalicFont=*-bolditalic,Scale=0.95]
% Sans para títulos, cuadros y figuras: TeX Gyre Heros.
\setsansfont{texgyreheros}[Extension=.otf,UprightFont=*-regular,BoldFont=*-bold,
  ItalicFont=*-italic,BoldItalicFont=*-bolditalic,Scale=MatchLowercase]
\setmonofont{Latin Modern Mono}[Scale=0.85]
\usepackage{polyglossia}
\setdefaultlanguage{spanish}
\setotherlanguage{english}
\usepackage[letterpaper,margin=27mm,headsep=7mm,footskip=11mm]{geometry}
\usepackage{microtype,graphicx,amsmath}
\usepackage{xcolor,placeins,etoolbox}
\usepackage{booktabs,tabularx,threeparttable,ragged2e}
\usepackage{enumitem}
\usepackage{titlesec}
\usepackage{caption}
\usepackage{fancyhdr}
\usepackage[hidelinks,unicode]{hyperref}
\usepackage{xurl}
\usepackage[autostyle]{csquotes}

% --- Paleta: un acento y un gris. Los cuadros permanecen en negro. -----------
\definecolor{tinta}{RGB}{23,59,54}    % acento (títulos, rótulos, trazos)
\definecolor{gris}{RGB}{107,107,102}  % cabeceras y notas discretas
\definecolor{humo}{RGB}{213,221,212}  % filetes muy suaves (solo figuras)

% --- Párrafo y página ----------------------------------------------------------
\linespread{1.04}
\setlength{\parskip}{0.3em}
\setlength{\parindent}{1.1em}
\setlength{\emergencystretch}{3em}
\tolerance=600 \hbadness=10000
\widowpenalty=10000 \clubpenalty=10000
\setlist{nosep}
\setcounter{secnumdepth}{3}

% --- Títulos de sección: sans, acento, numeración con punto (nivel 1) ---------
\titleformat{\section}{\sffamily\bfseries\large\color{tinta}}{\thesection.}{0.6em}{}
\titlespacing*{\section}{0pt}{1.6em}{0.6em}
\titleformat{\subsection}{\sffamily\bfseries\normalsize}{\thesubsection}{0.6em}{}
\titlespacing*{\subsection}{0pt}{1.1em}{0.35em}
\titleformat{\subsubsection}{\sffamily\itshape\normalsize}{\thesubsubsection}{0.6em}{}
\titleformat{\paragraph}[runin]{\itshape}{}{0pt}{}[.]

% --- Leyendas: rótulo en sans y acento, texto pequeño, alineado a la izquierda --
\captionsetup{font=small,labelfont={sf,bf,color=tinta},labelsep=period,
  justification=raggedright,singlelinecheck=false,skip=5pt}
\captionsetup[table]{position=above,skip=3pt}
\captionsetup[figure]{position=below}

% --- Cuadros: sans pequeña, sin líneas internas ni sombreado (guía del Boletín)
% Los flotantes reinician la fuente; el gancho correcto es \@floatboxreset.
\makeatletter
\g@addto@macro\@floatboxreset{\centering\ifdefstring{\@captype}{table}{\sffamily\small}{}}
\makeatother
\renewcommand{\arraystretch}{1.15}
\setlength{\tabcolsep}{5pt}
\setlength{\aboverulesep}{0.45ex}
\setlength{\belowrulesep}{0.65ex}
\newcolumntype{Y}{>{\RaggedRight\arraybackslash}X}  % texto, sin justificar
\newcolumntype{Z}{>{\RaggedLeft\arraybackslash}X}   % cifras alineadas a la derecha
\newcommand{\nd}{---}                               % dato no documentado
\newcommand{\aprox}{\ensuremath{\approx}}
\newcommand{\mayor}{\ensuremath{>}}
\newcommand{\menor}{\ensuremath{<}}
\newcommand{\cabecera}[1]{\textbf{#1}}              % cabecera de columna
\newcommand{\fila}[1]{\textbf{#1}}                  % rótulo de fila
% Notas al pie de cuadro: \begin{tablenotes}[flushleft]\footnotesize ... \end{tablenotes}
\appto\TPTnoteSettings{\footnotesize\color{gris}\RaggedRight}

% --- Cabecera y pie ------------------------------------------------------------
\pagestyle{fancy}\fancyhf{}
\renewcommand{\headrulewidth}{0pt}
\setlength{\headheight}{14pt}
\fancyhead[L]{\footnotesize\sffamily\color{gris}\papershorttitle}
\fancyhead[R]{\footnotesize\sffamily\color{gris}%
  \ifdefstring{\paperstatus}{draft}{Borrador \textperiodcentered{} fuentes: \sourcerevision}{}}
\fancyfoot[C]{\small\sffamily\thepage}
\fancypagestyle{plain}{\fancyhf{}\renewcommand{\headrulewidth}{0pt}%
  \fancyfoot[C]{\small\sffamily\thepage}%
  \fancyhead[R]{\footnotesize\sffamily\color{gris}%
    \ifdefstring{\paperstatus}{draft}{Borrador \textperiodcentered{} fuentes: \sourcerevision}{}}}

% --- Bloque de título (datos desde build/metadata.tex) -------------------------
\newcommand{\paperheading}{%
  \thispagestyle{plain}%
  \begin{center}
    {\sffamily\bfseries\Large\color{tinta}\papertitle\par}\vspace{0.55em}
    {\itshape\papertitleen\par}\vspace{1.2em}
    {\paperauthors\par}\vspace{0.45em}
    {\small\paperaffiliations\par}\vspace{0.25em}
    {\small\color{gris}Correspondencia: \papercontact\par}
  \end{center}
  \vspace{0.3em}{\color{tinta}\hrule height 0.5pt}\vspace{1.1em}}

% --- Resúmenes y palabras clave: medida algo más estrecha, rótulo en acento ---
\newenvironment{bloqueresumen}[1]{%
  \par\begingroup\small
  \setlength{\leftskip}{7mm}\setlength{\rightskip}{7mm}\setlength{\parindent}{0pt}
  \noindent{\sffamily\bfseries\color{tinta}#1.}\enspace\ignorespaces}{\par\endgroup}
\newcommand{\clave}[2]{\par\begingroup\small
  \setlength{\leftskip}{7mm}\setlength{\rightskip}{7mm}\setlength{\parindent}{0pt}
  \noindent{\sffamily\bfseries\color{tinta}#1:}\enspace #2\par\endgroup\medskip}

% --- Compatibilidad con el verificador: el banner ya no se invoca en el cuerpo --
\newcommand{\draftnotice}{}
\newcommand{\editorialtodo}[1]{\textbf{[Revisión editorial: #1]}}
